We say that a Hadamard manifold $M$ is \emph{asymptotically harmonic about a point $\xi\in M_\infty$}
if there is a number $h\ge0$ such that each horosphere in $M$ with center $\xi$ has constant mean curvature $h$.
Hyperbolic spaces $H^k_{\F}$, endowed with Fubini--Study metrics, are asymptotically harmonic about any point in $M_\infty$, where $\lambda_0=h^2/4>0$.
Asymptotically harmonic manifolds in the usual sense are asymptotically harmonic about any point in $M_\infty$.
Such manifolds arise, for example, in investigations on the rigidity of geodesic flows.
Furthermore, since horospheres are limits of spheres,
harmonic Hadamard manifolds are asymptotically harmonic about any point in $M_\infty$.
However, homogeneous Hadamard manifolds without Euclidean factor,
which are not hyperbolic spaces, are asymptotically harmonic about some, but not any point $M_\infty$.
A second main result of this article is

\begin{Theorem}\label{onlyacm}
If $K_M<0$ and $M$ is asymptotically harmonic about a point $\xi\in M_\infty$,
then the spectrum of $\Delta$ on $M$ is absolutely continuous with $\sigma(M)\subseteq\bigl[h^2/4,\infty\bigr)$,
where $h>0$ is the mean curvature of the horospheres with center $\xi$.
Moreover, if the sectional curvature of $M$ is negatively pinched and the covariant derivative of the curvature tensor of $M$ is uniformly bounded, then $\sigma(M)=\bigl[h^2/4,\infty\bigr)$.
\end{Theorem}

Say that a group of isometries of a Hadamard manifold $M$ is \emph{elementary}
if it is discrete and without torsion, fixes a point $\xi\in M_\infty$,
and such that it is of one of the following two types:
\begin{enumerate}[label=(\alph*)]\itemsep=0pt
\item\label{ax} $\Gamma$ leaves a geodesic $\gamma$ in $M$ emanating from $\xi$ invariant.
\item\label{pa} $\Gamma$ leaves Busemann functions associated to $\xi$ invariant.
\end{enumerate}
In the first case, $C=\Gamma\backslash\gamma$ is a circle and the projection $M\to\gamma$
along horospheres with center~$\xi$ descends to a Riemannian submersion $\pi\colon N\to C$,
where $N=\Gamma\backslash M$.
In the second case,
any Busemann function $b$ associated to $\xi$ is invariant under $\Gamma$
and hence pushes down to $N$ and defines a Riemannian submersion $\pi=b\colon N\to\R$.

\begin{Theorem}\label{onlyacn}
For $M$ as in Theorem~{\rm \ref{onlyacm}}, if $\Gamma$ is an elementary group of isometries of $M$ fixing~$\xi$,
then the spectrum of the quotient $N=\Gamma\backslash M$ is absolutely continuous
with $\sigma(N)\subseteq\bigl[h^2/4,\infty\bigr)$, where $h>0$ denotes the mean curvature of the fibers of $\pi$.
Moreover,
\begin{enumerate}[label={\rm (\arabic*)}]\itemsep=0pt
\item
if $N$ is of type {\rm \ref{pa}} and the fibers of $\pi$ are of finite volume or
\item
if $N$ is of type {\rm \ref{pa}}, the fibers of $\pi$ are of infinite volume,
the sectional curvature of $M$ is negatively pinched,
and the covariant derivative of the curvature tensor of $M$ is uniformly bounded,
\end{enumerate}
then $\sigma(N)=\bigl[h^2/4,\infty\bigr)$.
\end{Theorem}
